
\documentclass[uplatex,dvipdfmx,a4paper,10pt]{jsarticle}
\usepackage[margin=23mm]{geometry}
\usepackage{longtable,booktabs,array}
\usepackage{enumitem}
\usepackage{xcolor}
\usepackage{tcolorbox}
\usepackage{hyperref}
\hypersetup{colorlinks=true,linkcolor=black,urlcolor=black}
\setlength{\parindent}{0pt}
\setlength{\parskip}{0.65em}
\renewcommand{\baselinestretch}{1.08}
\newtcolorbox{infobox}[1]{colback=gray!8,colframe=gray!55,title=#1,fonttitle=\bfseries,arc=1mm,boxrule=0.4pt}
\newtcolorbox{todobox}{colback=gray!5,colframe=black!50,title=TODO（図）と画像生成用プロンプト,fonttitle=\bfseries,arc=1mm,boxrule=0.4pt}
\begin{document}
\begin{titlepage}
\centering
{\Large 17 章 数学的基礎\par}
\vspace{0.8em}
{\large Mathematical Foundations の日本語まとめ\par}
\vspace{0.8em}
{SWEBOK Guide v4.0a Chapter 17 を初学者向けに再構成\par}
\vspace{2em}
{作成日：2026 5 6\par}
\vfill
{専門用語は原則として日本語に直し、必要に応じて英語を括弧内に併記した。文体は「だ・である調」に統一した。\par}
\end{titlepage}
\tableofcontents
\newpage
\section{全体概要：この章で学ぶこと}
第17章は、ソフトウェア技術者が「曖昧でない考え方」をコード、仕様、モデル、テストへ落とし込むための数学を扱う。ここでいう数学は、計算練習としての算術だけではない。中心にあるのは、論理、集合、関係、関数、グラフ、状態、文法、確率、誤差、代数構造といった、ソフトウェアを正確に考えるための道具である。

この章の狙いは、複雑な対象を形式化し、推論し、妥当な結論を得るための基礎を理解することにある。たとえば、if文の条件、データ構造、状態遷移、パーサ、暗号、乱択アルゴリズム、数値シミュレーションは、いずれも数学的な概念に支えられている。

\begin{infobox}{この資料の主張}
数学的基礎は、ソフトウェア開発における「正しく考えるための共通語」である。論理は条件分岐と仕様の意味を支え、集合と関数はデータと処理の関係を支え、グラフと木は構造を支え、有限状態機械と文法は振る舞いと言語処理を支える。確率と誤差は、不確実性と近似を扱うために必要である。
\end{infobox}
到達目標 この章を読み終えると、次のことを説明できるようになる。

\begin{itemize}[leftmargin=2zw]
\item 命題論理と述語論理の違いを説明できる。
\item 直接証明、背理法、数学的帰納法、例による存在証明を区別できる。
\item 集合、関係、関数をデータの整理方法として説明できる。
\item グラフ、木、二分探索木、有限状態機械がソフトウェアで使われる理由を説明できる。
\item 文法、正規表現、文脈自由文法がプログラミング言語や入力処理と関係する理由を説明できる。
\item 数え上げ、離散確率、数値誤差を、テスト、性能評価、シミュレーションに結び付けて説明できる。
\end{itemize}
学習順序 初学者は、まず論理と集合を読むとよい。これらは後のグラフ、木、状態機械、文法、確率の土台になる。次に、データ構造に近いグラフと木、振る舞いに近い有限状態機械と文法を読む。最後に、数論、数え上げ、確率、数値誤差、代数構造、微積分を、用途と結び付けて読むとよい。

\begin{todobox}
17章の全体地図を入れる。

画像生成用プロンプト

白背景、モノクロ、A4本文幅に収まる横長の学習用図解を作成する。中央に「数学的基礎」を置き、周囲に「論理」「証明」「集合・関係・関数」「グラフ・木」「有限状態機械」「文法」「数論」「数え上げ」「離散確率」「数値誤差」「代数構造」「工学微積分」「新しい応用」を配置する。矢印で「正確な記述」「構造の理解」「振る舞いの理解」「不確実性と近似の扱い」を示す。日本語ラベル、講義ノート風、細線、余白広め。
\end{todobox}
\subsection{最初に必要な用語}
\begin{longtable}{>{\raggedright\arraybackslash}p{0.26\linewidth} >{\raggedright\arraybackslash}p{0.68\linewidth}}
\toprule
用語 & 初学者向けの説明 \\
\midrule
\endfirsthead
\toprule
用語 & 初学者向けの説明 \\
\midrule
\endhead
形式的 & 意味が曖昧にならないよう、記号、規則、推論手順を明確に定めること。 \\
命題 & 真または偽のどちらかに判定できる文である。 \\
述語 & 対象について性質や関係を述べる文の型である。「xは赤い」「xはyより大きい」などである。 \\
証明 & ある主張が必ず真であることを、認めた規則から厳密に示す説明である。 \\
集合 & 対象をひとまとまりにしたものである。要素の順序や重複ではなく、含まれるかどうかに注目する。 \\
関係 & 複数の集合の要素どうしの対応である。例として、人と電話番号の対応がある。 \\
関数 & 入力ごとに出力がただ一つ決まる関係である。 \\
グラフ & 点と線で対象同士のつながりを表す構造である。 \\
木 & 根から枝分かれする階層構造であり、単純な閉路を持たない。 \\
有限状態機械 & 有限個の状態を持ち、入力によって状態を遷移する抽象モデルである。 \\
文法 & 文字列や文を作る規則である。プログラミング言語や設定ファイルの構文を扱う基礎である。 \\
離散確率 & 数え上げられる結果に確率を割り当てる考え方である。 \\
数値誤差 & 計算機で表した近似値と本当の値との差である。 \\
\bottomrule
\end{longtable}
\section{導入：なぜ数学的基礎がソフトウェアに必要なのか}
ソフトウェアは、曖昧なままでは実行できない。人間同士の会話では「だいたい」「普通は」「適切に」といった表現が許されることがある。しかし、コード、仕様、テスト、モデルでは、どの条件でどの処理を行うかを明確にする必要がある。数学的基礎は、この明確化を支える。

数学は、数だけを扱う学問ではない。より広く見れば、数学は形式体系を扱う。形式体系とは、対象、記号、規則、推論方法を定め、その中で何が成り立つかを調べる枠組みである。ソフトウェア技術者は、業務、通信、画面操作、状態遷移、データ、セキュリティ規則などを形式化しなければならない。

\begin{infobox}{重要}
第17章の数学は、計算式を速く解くためだけの数学ではない。ソフトウェアで「何が正しいか」「どの入力でどの出力になるか」「どの状態からどの状態へ進むか」「どの文字列が有効か」を曖昧なく扱うための数学である。
\end{infobox}
\section{1 基本論理}
基本論理は、真偽を扱うための土台である。仕様の条件、プログラムの分岐、テストの期待結果、セキュリティ規則は、多くの場合、論理式として読める。論理を理解すると、「同じ意味の条件」「矛盾する条件」「常に真になる条件」を見分けやすくなる。

\subsection{1.1 命題論理}
命題とは、真または偽のどちらかであり、両方ではない文である。「太陽は恒星である」「2 + 3 = 5」は命題である。一方、「a + 3 = b」は、aとbの値が決まらない限り真偽を判定できないため、そのままでは命題ではない。

命題論理は、命題を論理演算で組み合わせる方法である。基本演算には、否定、論理積、論理和、排他的論理和、含意がある。プログラムでは、これらは not、and、or、xor、条件付きの意味として現れる。

\begin{longtable}{>{\raggedright\arraybackslash}p{0.26\linewidth} >{\raggedright\arraybackslash}p{0.68\linewidth}}
\toprule
考え方 & 説明 \\
\midrule
\endfirsthead
\toprule
考え方 & 説明 \\
\midrule
\endhead
排中律 & 任意の命題pについて、pは真であるか偽であるかのどちらかである。 \\
矛盾律 & 任意の命題pについて、pが真かつ偽であることはない。 \\
恒真式 & どの真理値を入れても常に真になる複合命題である。 \\
矛盾式 & どの真理値を入れても常に偽になる複合命題である。 \\
論理的同値 & 二つの論理式が、常に同じ真理値を持つ関係である。 \\
ド・モルガンの法則 & 「AかつBではない」は「Aでない、またはBでない」に等しい、という形の変形規則である。 \\
\bottomrule
\end{longtable}
\begin{infobox}{ソフトウェアでの見方}
複雑なif条件は、命題論理として整理できる。たとえば「管理者であり、かつ停止中でない場合だけ操作できる」という条件は、権限判定の論理式である。ド・モルガンの法則を使うと、否定条件を読みやすく変形できる。
\end{infobox}
\begin{todobox}
命題論理の真理値表と論理演算の図を入れる。

画像生成用プロンプト

白背景、モノクロ、A4本文幅の表形式図を作成する。左にpとqの真偽4通り、右に「否定」「論理積」「論理和」「含意」の結果を並べる。下部に「条件分岐、アクセス制御、テスト期待結果に使う」と注記する。日本語ラベル、講義ノート風、細線。
\end{todobox}
\subsection{1.2 述語論理}
述語論理は、命題論理を拡張し、変数を含む性質や関係を扱えるようにした論理である。たとえば「xは赤い」「xは利用者である」「xはyに依存している」といった文の型を述語と呼ぶ。xやyに具体的な対象を入れると、真偽を判定できる文になる。

述語論理では、量化子が重要である。全称量化子は「すべてのxについて」を表し、存在量化子は「少なくとも一つのxが存在する」を表す。仕様では、「すべての注文は一つの顧客に属する」「少なくとも一人の管理者が存在する」といった制約を表すために使える。

\begin{longtable}{>{\raggedright\arraybackslash}p{0.26\linewidth} >{\raggedright\arraybackslash}p{0.68\linewidth}}
\toprule
量化子 & 意味と例 \\
\midrule
\endfirsthead
\toprule
量化子 & 意味と例 \\
\midrule
\endhead
全称量化子 & 「すべてのxについて成り立つ」を表す。例：すべての認証済み利用者は利用者IDを持つ。 \\
存在量化子 & 「少なくとも一つのxについて成り立つ」を表す。例：少なくとも一つの管理者アカウントが存在する。 \\
束縛変数 & 量化子の範囲内で意味が決まる変数である。 \\
自由変数 & 量化子に束縛されておらず、値が外から与えられる変数である。 \\
\bottomrule
\end{longtable}
\begin{infobox}{命題論理との違い}
命題論理は「この文は真か偽か」を扱う。述語論理は「どの対象について成り立つか」「すべてに成り立つか」「少なくとも一つに成り立つか」を扱う。データベース制約、クラス不変条件、業務規則の表現では、述語論理の考え方が特に重要である。
\end{infobox}
\begin{todobox}
量化子の意味を示す図を入れる。

画像生成用プロンプト

白背景、モノクロ、A4本文幅の概念図を作成する。集合「利用者」を円で表し、全称量化子は全員にチェック印、存在量化子は少なくとも一人にチェック印を付ける。例として「全利用者はIDを持つ」「少なくとも一人は管理者である」を添える。日本語ラベル、講義ノート風。
\end{todobox}
\section{2 証明技法}
証明とは、ある主張が真であることを、仮定、定義、すでに証明された命題から厳密に導くことである。ソフトウェア開発では、証明という語を日常的に使わなくても、設計の正しさ、アルゴリズムの停止性、データ構造の性質、テストケースの十分性を考えるときに証明の考え方が必要になる。

\begin{longtable}{>{\raggedright\arraybackslash}p{0.26\linewidth} >{\raggedright\arraybackslash}p{0.68\linewidth}}
\toprule
用語 & 説明 \\
\midrule
\endfirsthead
\toprule
用語 & 説明 \\
\midrule
\endhead
公理・仮定 & 証明の出発点として受け入れる前提である。 \\
定理 & 証明によって真であると示された主張である。 \\
補題 & 別の定理を証明するために使う小さな定理である。 \\
系 & 証明済みの定理から直接導かれる主張である。 \\
予想 & 真偽がまだ確定していない主張である。 \\
\bottomrule
\end{longtable}
\subsection{2.1 直接証明}
直接証明は、「pならばq」を示すとき、pが真であると仮定し、定義や計算を使ってqが真であることを導く方法である。最も素直な証明方法であり、アルゴリズムの性質を説明するときにも使いやすい。

\begin{infobox}{例}
「nが奇数なら、nの2乗から1を引いた数は偶数である」を示すとき、n = 2k + 1 と置く。すると n\textasciicircum{}2 は 4k\textasciicircum{}2 + 4k + 1 となり、n\textasciicircum{}2 - 1 は 4k\textasciicircum{}2 + 4k である。これは2で割り切れるため偶数である。
\end{infobox}
\subsection{2.2 背理法と対偶証明}
背理法は、示したい主張が偽であると仮定し、その仮定から矛盾を導く方法である。矛盾が出るなら、最初の仮定が誤りであり、示したい主張は真であると結論づける。

対偶証明は、「pならばq」を示す代わりに、「qでないならpでない」を示す方法である。論理的には、この二つは同値である。直接証明が難しいときに有効である。

\subsection{2.3 数学的帰納法}
数学的帰納法は、自然数で番号付けされた無限個の主張を証明するための技法である。第一に、最初の場合が成り立つことを示す。第二に、ある任意のkで成り立つと仮定したとき、k+1でも成り立つことを示す。この二つにより、すべての自然数について成り立つと結論できる。

\begin{longtable}{>{\raggedright\arraybackslash}p{0.26\linewidth} >{\raggedright\arraybackslash}p{0.68\linewidth}}
\toprule
段階 & 意味 \\
\midrule
\endfirsthead
\toprule
段階 & 意味 \\
\midrule
\endhead
基底段階 & 最初のケース、通常はn = 1またはn = 0で主張が成り立つことを示す。 \\
帰納法の仮定 & 任意のkで主張が成り立つと仮定する。 \\
帰納段階 & kで成り立つならk+1でも成り立つことを示す。 \\
\bottomrule
\end{longtable}
\begin{infobox}{ソフトウェアでの見方}
再帰関数、ループ不変条件、木構造の性質、リスト処理の正しさは、数学的帰納法と相性がよい。「空リストで正しい」「先頭を一つ処理して残りで正しいなら全体で正しい」という考え方は帰納法そのものである。
\end{infobox}
\subsection{2.4 例による証明}
例による証明が有効なのは、「存在する」ことを示したい場合である。たとえば「条件を満たす入力が少なくとも一つある」を示すには、その入力例を一つ示せばよい。

一方で、「すべての場合に成り立つ」ことは、いくつかの例だけでは証明できない。例を多く見せても、一般の証明にはならない。これは不適切な一般化と呼べる。テストで多数のケースが通っても、欠陥が存在しないことの証明にはならないという点にもつながる。

\begin{todobox}
証明技法の比較フローを入れる。

画像生成用プロンプト

白背景、モノクロ、A4本文幅の判定フローを作成する。開始は「証明したい主張」。分岐は「存在を示すか」「pならばqを示すか」「自然数に関する主張か」「直接示しにくいか」。到達先に「例による存在証明」「直接証明」「数学的帰納法」「背理法・対偶証明」を配置する。日本語ラベル、講義ノート風。
\end{todobox}
\section{3 集合・関係・関数}
集合、関係、関数は、データと処理を正確に考えるための基本語彙である。データベースの表、APIの入力と出力、型、制約、関連、写像は、集合・関係・関数の考え方で整理できる。

\subsection{3.1 集合の基本}
集合とは、対象の集まりである。対象は要素と呼ばれる。集合では、同じ要素が含まれるかどうかが重要であり、通常は並び順や重複は本質ではない。

\begin{longtable}{>{\raggedright\arraybackslash}p{0.26\linewidth} >{\raggedright\arraybackslash}p{0.68\linewidth}}
\toprule
概念 & 説明 \\
\midrule
\endfirsthead
\toprule
概念 & 説明 \\
\midrule
\endhead
有限集合 & 要素数が有限である集合である。 \\
無限集合 & 要素数が有限ではない集合である。自然数全体などである。 \\
濃度 & 有限集合では要素数を意味する。 \\
全体集合 & 議論している範囲のすべての要素を含む集合である。 \\
部分集合 & 集合Xのすべての要素が集合Yにも含まれるとき、XはYの部分集合である。 \\
真部分集合 & 部分集合であり、かつ等しくない集合である。 \\
空集合 & 要素を一つも持たない集合である。 \\
冪集合 & ある集合のすべての部分集合を集めた集合である。 \\
\bottomrule
\end{longtable}
\subsection{3.2 集合演算}
集合演算は、複数の集合を組み合わせたり、差を取ったりする操作である。検索条件、権限判定、データ抽出、テスト対象の分類で頻繁に現れる。

\begin{longtable}{>{\raggedright\arraybackslash}p{0.26\linewidth} >{\raggedright\arraybackslash}p{0.68\linewidth}}
\toprule
演算 & 意味とソフトウェアでの例 \\
\midrule
\endfirsthead
\toprule
演算 & 意味とソフトウェアでの例 \\
\midrule
\endhead
共通部分 & 両方の集合に含まれる要素である。例：有料会員かつ管理者である利用者。 \\
和集合 & どちらか一方、または両方に含まれる要素である。例：メール通知対象またはSMS通知対象の利用者。 \\
補集合 & 全体集合のうち、指定集合に含まれない要素である。例：退会済みでない利用者。 \\
差集合 & 一方に含まれ、他方に含まれない要素である。例：登録済みだが本人確認未完了の利用者。 \\
直積 & 二つの集合から順序付き対を作る演算である。例：利用者集合と権限集合の組合せ。 \\
\bottomrule
\end{longtable}
\begin{todobox}
集合演算のベン図を入れる。

画像生成用プロンプト

白背景、モノクロ、A4本文幅の4分割図を作成する。各区画に「共通部分」「和集合」「補集合」「差集合」を置き、2つの円と全体集合の長方形で塗り分ける。各図に短いソフトウェア例を添える。日本語ラベル、講義ノート風、印刷で見やすい。
\end{todobox}
\subsection{3.3 集合の性質}
集合には、結合法則、交換法則、分配法則、恒等法則、補法則、冪等法則、吸収法則、ド・モルガンの法則などがある。これらは条件式や検索条件を変形するときに役立つ。

\begin{infobox}{実務での例}
「Aに属し、かつBまたはCに属する」を「AかつB、またはAかつC」と変形すると、検索条件や権限条件を実装しやすくなる場合がある。これは集合の分配法則に対応する。
\end{infobox}
\subsection{3.4 関係と関数}
関係とは、集合の要素どうしの対応である。たとえば、人と電話番号の対応、注文と商品明細の対応、モジュールと依存先の対応は関係である。

関数とは、入力ごとに出力がただ一つ決まる関係である。入力が同じなのに複数の異なる出力があり得るなら、それは関数ではない。関数は、プログラムの処理、API、写像、変換、計算式を考える基本である。

\begin{longtable}{>{\raggedright\arraybackslash}p{0.26\linewidth} >{\raggedright\arraybackslash}p{0.68\linewidth}}
\toprule
用語 & 説明 \\
\midrule
\endfirsthead
\toprule
用語 & 説明 \\
\midrule
\endhead
定義域 & 入力として取り得る値の集合である。 \\
値域 & 対応する出力値の集合である。 \\
関係 & 入力と出力の対応である。入力一つに複数出力が対応してもよい。 \\
関数 & 入力一つに対して出力が一つだけ決まる関係である。 \\
垂直線テスト & グラフ上で同じxに複数のyが対応するなら関数ではない、という確認方法である。 \\
\bottomrule
\end{longtable}
\begin{infobox}{API設計での見方}
同じリクエストを同じ状態で送ったときに同じレスポンスが返るなら、APIは関数のように扱いやすい。一方、時刻、乱数、外部状態、ユーザセッションによって出力が変わる場合は、状態も入力として考えなければならない。
\end{infobox}
\begin{todobox}
関係と関数の違いを示す図を入れる。

画像生成用プロンプト

白背景、モノクロ、A4本文幅の左右比較図を作成する。左は「関係」として一つの入力から複数の出力へ矢印を出す。右は「関数」として各入力からただ一つの出力へ矢印を出す。下部に「同じ入力に複数出力があると関数ではない」と注記する。日本語ラベル、講義ノート風。
\end{todobox}
\section{4 グラフと木}
グラフと木は、つながりと階層を表す構造である。ネットワーク、依存関係、状態遷移、制御フロー、データベースの関連、ファイルシステム、構文木などは、グラフまたは木として考えられる。

\subsection{4.1 グラフ}
グラフは、頂点と辺からなる構造である。頂点は対象を表し、辺は対象間のつながりを表す。辺は、有向でも無向でもよく、重みを持つこともある。

\begin{longtable}{>{\raggedright\arraybackslash}p{0.26\linewidth} >{\raggedright\arraybackslash}p{0.68\linewidth}}
\toprule
グラフの種類 & 説明と例 \\
\midrule
\endfirsthead
\toprule
グラフの種類 & 説明と例 \\
\midrule
\endhead
単純グラフ & 同じ頂点対を結ぶ辺が一つまでで、自己ループを持たないグラフである。 \\
多重グラフ & 同じ頂点対の間に複数の辺を許すグラフである。複数種類の関係を表す場合に使える。 \\
擬グラフ & 頂点から同じ頂点へ戻る自己ループを許すグラフである。 \\
有向グラフ & 辺に向きがあるグラフである。依存関係、遷移、呼び出し関係に適する。 \\
重み付きグラフ & 辺に距離、費用、遅延、容量などの数値を持たせたグラフである。 \\
\bottomrule
\end{longtable}
グラフでは、隣接、接続、端点、次数、入次数、出次数、経路、回路、閉路などの用語を使う。これらは、ネットワーク経路探索、依存関係解析、ビルド順序決定、到達可能性分析で重要である。

\begin{longtable}{>{\raggedright\arraybackslash}p{0.26\linewidth} >{\raggedright\arraybackslash}p{0.68\linewidth}}
\toprule
表現方法 & 特徴 \\
\midrule
\endfirsthead
\toprule
表現方法 & 特徴 \\
\midrule
\endhead
隣接リスト & 各頂点について隣接する頂点を列挙する。辺が少ない疎なグラフで効率がよい。 \\
隣接行列 & 頂点対ごとに辺の有無を行列で表す。辺の有無を高速に確認しやすい。 \\
接続行列 & 頂点と辺の関係を行列で表す。グラフ理論の解析で使われる。 \\
\bottomrule
\end{longtable}
\begin{todobox}
グラフの種類を比較する図を入れる。

画像生成用プロンプト

白背景、モノクロ、A4本文幅の横長比較図を作成する。5つの小図を並べ、「単純グラフ」「多重グラフ」「擬グラフ」「有向グラフ」「重み付きグラフ」とラベルを付ける。各小図は3つ程度の頂点で表す。重み付きグラフは辺に数値を付ける。日本語ラベル、講義ノート風。
\end{todobox}
\subsection{4.2 木}
木は、根を起点に親子関係で枝分かれする階層構造である。別の見方をすると、木は連結で単純な閉路を持たない無向グラフである。木の辺数は、節点数から1を引いた数になる。

\begin{longtable}{>{\raggedright\arraybackslash}p{0.26\linewidth} >{\raggedright\arraybackslash}p{0.68\linewidth}}
\toprule
用語 & 説明 \\
\midrule
\endfirsthead
\toprule
用語 & 説明 \\
\midrule
\endhead
根 & 木の出発点となる節点である。 \\
親と子 & ある節点から下位の節点へ向かう関係である。根以外の節点は一つの親を持つ。 \\
内部節点 & 少なくとも一つの子を持つ節点である。 \\
葉 & 子を持たない節点である。 \\
レベル & 根から何本の辺をたどるかで決まる深さである。 \\
高さ & 木の中で最大のレベルである。 \\
兄弟 & 同じ親を持つ節点である。 \\
\bottomrule
\end{longtable}
二分木は、各節点が高々二つの子を持つ木である。二分探索木は、各節点にキーを持たせ、左部分木には小さいキー、右部分木には大きいキーを置くようにした二分木である。探索、挿入、整列に使われる。

\begin{longtable}{>{\raggedright\arraybackslash}p{0.26\linewidth} >{\raggedright\arraybackslash}p{0.68\linewidth}}
\toprule
二分木の種類 & 説明 \\
\midrule
\endfirsthead
\toprule
二分木の種類 & 説明 \\
\midrule
\endhead
二分木 & 各内部節点の子が高々2個であり、左と右の位置に意味がある木である。 \\
満二分木 & すべての内部節点がちょうど2個の子を持つ二分木である。 \\
完全二分木 & 最後のレベル以外が埋まり、最後のレベルは左から詰まっている二分木である。 \\
平衡二分木 & 葉の深さが大きく偏らないよう保たれた二分木である。 \\
二分探索木 & 左は小さい値、右は大きい値という順序規則を持つ二分木である。 \\
\bottomrule
\end{longtable}
木の巡回には、前順、中順、後順がある。前順は根を先に訪れる。中順は左部分木、根、右部分木の順に訪れる。後順は部分木を先に訪れ、根を最後に訪れる。二分探索木を中順に巡回すると、キーを昇順に取り出せる。

\begin{todobox}
木構造と二分探索木の巡回を示す図を入れる。

画像生成用プロンプト

白背景、モノクロ、A4本文幅の図を作成する。左に一般の木を描き、根、親、子、葉、レベル、高さをラベル付けする。右に二分探索木を描き、中順巡回で値が昇順になることを矢印で示す。日本語ラベル、講義ノート風。
\end{todobox}
\section{5 有限状態機械}
有限状態機械は、有限個の状態と、入力による状態遷移で振る舞いを表す抽象モデルである。画面遷移、通信プロトコル、ワークフロー、字句解析器、組込み制御、ゲームの状態管理などに使える。

状態機械では、システムはある状態にいる。入力が来ると、遷移関数に従って別の状態へ移り、必要なら出力を出す。開始状態から始まり、受理状態や終了状態に到達することで、処理の成功を表せる。

\begin{longtable}{>{\raggedright\arraybackslash}p{0.26\linewidth} >{\raggedright\arraybackslash}p{0.68\linewidth}}
\toprule
構成要素 & 説明 \\
\midrule
\endfirsthead
\toprule
構成要素 & 説明 \\
\midrule
\endhead
状態集合 & 取り得る状態の有限集合である。 \\
入力記号集合 & 状態遷移を引き起こす入力の集合である。 \\
出力記号集合 & 遷移時または状態で出される出力の集合である。 \\
遷移関数 & 現在状態と入力から次状態を決める規則である。 \\
出力関数 & 現在状態と入力から出力を決める規則である。 \\
初期状態 & 処理を開始する状態である。 \\
受理状態 & 入力列が有効に処理されたことを示す状態である。 \\
\bottomrule
\end{longtable}
\begin{infobox}{実務での見方}
状態機械を使うと、「この状態でこの入力が来たらどうするか」を表にできる。漏れた入力、到達不能な状態、不正な遷移を見つけやすくなる。これはテストケース設計にも有用である。
\end{infobox}
\begin{todobox}
有限状態機械の状態遷移図と状態表を入れる。

画像生成用プロンプト

白背景、モノクロ、A4本文幅の左右比較図を作成する。左に3状態S0、S1、S2の有向状態遷移図を描き、入力0/1で遷移する矢印を付ける。右に「現在状態」「入力」「次状態」「出力」の状態表を置く。下部に「図と表は同じ情報を表す」と注記する。日本語ラベル、講義ノート風。
\end{todobox}
\section{6 文法}
文法は、文字列や文を作る規則である。自然言語では文法が「文として自然か」を判断するが、形式言語では文法が「その文字列が言語に属するか」を厳密に決める。プログラミング言語、正規表現、設定ファイル、問い合わせ言語、データ形式は、形式文法と深く関係している。

\subsection{6.1 形式言語と文法}
形式言語は、有限のアルファベット上の有限長文字列の集合である。文法は、その言語に属する文字列を生成する規則を与える。すべての正しい文字列を列挙するのではなく、生成規則によってコンパクトに定義する点が重要である。

\begin{longtable}{>{\raggedright\arraybackslash}p{0.26\linewidth} >{\raggedright\arraybackslash}p{0.68\linewidth}}
\toprule
用語 & 説明 \\
\midrule
\endfirsthead
\toprule
用語 & 説明 \\
\midrule
\endhead
アルファベット & 文字列を構成する基本記号の集合である。 \\
終端記号 & 最終的な文字列に現れる記号である。 \\
非終端記号 & 生成途中の概念を表す記号である。式、項、文などに相当する。 \\
開始記号 & 生成を開始する非終端記号である。 \\
生成規則 & ある記号列を別の記号列に置き換える規則である。 \\
文法が生成する言語 & 開始記号から生成規則を適用して得られるすべての終端記号列の集合である。 \\
\bottomrule
\end{longtable}
\subsection{6.2 言語認識とチョムスキー階層}
形式文法は、許す生成規則の種類によって分類できる。代表的な分類がチョムスキー階層である。正規文法は最も制約が強く、句構造文法は最も一般的である。

\begin{longtable}{>{\raggedright\arraybackslash}p{0.26\linewidth} >{\raggedright\arraybackslash}p{0.68\linewidth}}
\toprule
種類 & 説明と用途 \\
\midrule
\endfirsthead
\toprule
種類 & 説明と用途 \\
\midrule
\endhead
句構造文法 & 最も一般的な形式文法である。 \\
文脈依存文法 & 置換規則が周囲の文脈に依存し得る文法である。 \\
文脈自由文法 & 左辺が一つの非終端記号である文法である。多くのプログラミング言語の構文の基礎である。 \\
正規文法 & 正規表現や有限状態機械と関係が深い文法である。字句解析や簡単なパターン照合に使われる。 \\
\bottomrule
\end{longtable}
正規表現は、文字列のパターンを表す記法である。たとえば、ある文字列が特定の形式のID、メールアドレス、トークン、ログ行に合うかを検査するために使われる。ただし、正規表現は万能ではない。入れ子構造を持つプログラミング言語の構文全体を扱うには、文脈自由文法などが必要になる。

\begin{todobox}
文法とチョムスキー階層の図を入れる。

画像生成用プロンプト

白背景、モノクロ、A4本文幅の入れ子図を作成する。外側から「句構造文法」「文脈依存文法」「文脈自由文法」「正規文法」を入れ子の長方形で表す。右側に用途例として「一般的な形式言語」「自然言語モデル」「プログラミング言語構文」「正規表現・字句解析」を並べる。日本語ラベル、講義ノート風。
\end{todobox}
\section{7 数論}
数論は、整数を中心に数の性質を扱う分野である。ソフトウェアでは、暗号、ハッシュ、乱数、剰余計算、符号化、アルゴリズム解析で重要になる。

\subsection{7.1 数の種類}
\begin{longtable}{>{\raggedright\arraybackslash}p{0.26\linewidth} >{\raggedright\arraybackslash}p{0.68\linewidth}}
\toprule
数の種類 & 説明 \\
\midrule
\endfirsthead
\toprule
数の種類 & 説明 \\
\midrule
\endhead
自然数 & 1, 2, 3, ... のように数えるための数である。ただし0を含めるかは流儀により異なる。 \\
非負整数 & 0と自然数を合わせた集合である。 \\
整数 & 負の数、0、正の数を含む集合である。 \\
有理数 & 二つの整数の比として表せる数である。小数展開は有限または循環する。 \\
無理数 & 整数の比で表せない数である。小数展開は有限でも循環でもない。 \\
実数 & 有理数と無理数を合わせた数である。 \\
虚数 & iを含む数である。iは2乗すると-1になる数として扱われる。 \\
複素数 & 実部と虚部を持つa + biの形の数である。 \\
\bottomrule
\end{longtable}
\begin{todobox}
数の種類の包含図を入れる。

画像生成用プロンプト

白背景、モノクロ、A4本文幅の入れ子図を作成する。外側に「複素数」、内側に「実数」、実数の中に「有理数」「無理数」、有理数の中に「整数」、整数の中に「非負整数」「自然数」を配置する。虚数は実数と異なる方向に示す。日本語ラベル、講義ノート風。
\end{todobox}
\subsection{7.2 整除性・剰余・合同}
整数aが整数bを割り切るとは、b = a × cとなる整数cが存在することである。これを整除性という。剰余は、割り算で余る値である。合同は、二つの数の差がある整数mで割り切れる関係であり、剰余計算を扱う基礎である。

\begin{infobox}{ソフトウェアでの見方}
剰余は、配列の循環、ハッシュテーブル、時刻計算、暗号、チェックディジットに現れる。たとえば「曜日」は7で割った余りとして扱える。
\end{infobox}
\subsection{7.3 素数と最大公約数}
素数は、1と自分自身以外の正の約数を持たない1より大きい整数である。素数でない1より大きい整数は合成数である。大きな素数は公開鍵暗号などで重要である。

最大公約数は、二つの整数をともに割り切る最大の整数である。最大公約数が1である二つの整数は互いに素である。互いに素という性質は、剰余演算や暗号アルゴリズムで重要である。

\section{8 数え上げの基礎}
数え上げは、可能な選択肢や組合せの数を数えるための考え方である。テストケース数、入力の組合せ、経路数、状態数、パスワード探索空間、スケジューリング候補を見積もるときに使う。

\begin{longtable}{>{\raggedright\arraybackslash}p{0.26\linewidth} >{\raggedright\arraybackslash}p{0.68\linewidth}}
\toprule
規則 & 説明と例 \\
\midrule
\endfirsthead
\toprule
規則 & 説明と例 \\
\midrule
\endhead
和の規則 & 同時に起こらない選択肢AとBがあるとき、全体の数はAの数とBの数の和である。 \\
積の規則 & 一つ目の選択の後に二つ目の選択を行うとき、全体の数はそれぞれの数の積である。 \\
包除原理 & 重なりがある二つの集合を数えるとき、重複して数えた共通部分を引く。 \\
再帰 & 対象や処理を、自分自身のより小さい場合で定義する。 \\
順列 & 順序が意味を持つ選び方である。 \\
組合せ & 順序が意味を持たない選び方である。 \\
\bottomrule
\end{longtable}
\begin{infobox}{テスト設計での見方}
入力項目が増えると、すべての組合せを試す数は急増する。数え上げを理解すると、全探索が現実的か、同値分割や組合せテストで減らすべきかを判断しやすくなる。
\end{infobox}
\begin{todobox}
和の規則・積の規則・順列・組合せの比較図を入れる。

画像生成用プロンプト

白背景、モノクロ、A4本文幅の比較表風図を作成する。列に「和の規則」「積の規則」「順列」「組合せ」を置き、行に「何を数えるか」「順序は重要か」「代表例」を置く。下部に「テストケース爆発を見積もる道具」と注記する。日本語ラベル、講義ノート風。
\end{todobox}
\section{9 離散確率}
確率は、偶然性を数学的に扱うための言葉である。離散確率は、サイコロの目、テストの合否、障害発生有無、リクエスト種別など、数え上げられる結果に確率を割り当てる。

\begin{longtable}{>{\raggedright\arraybackslash}p{0.26\linewidth} >{\raggedright\arraybackslash}p{0.68\linewidth}}
\toprule
用語 & 説明 \\
\midrule
\endfirsthead
\toprule
用語 & 説明 \\
\midrule
\endhead
標本空間 & 起こり得るすべての結果の集合である。 \\
事象 & 標本空間の部分集合である。 \\
確率変数 & 結果に数値を割り当てる規則である。 \\
離散確率変数 & 取り得る値を列挙できる確率変数である。 \\
連続確率変数 & 値をすべて列挙できない範囲で扱う確率変数である。 \\
確率分布 & 値とその値が起こる確率の対応である。 \\
平均 & 長期的に期待される中心的な値である。 \\
分散・標準偏差 & 値のばらつきの大きさを表す。 \\
\bottomrule
\end{longtable}
離散確率分布では、各結果の確率は0以上1以下であり、すべての確率の総和は1である。たとえば、公平なサイコロでは、1から6までの各目の確率はすべて1/6である。

\begin{infobox}{ソフトウェアでの見方}
ランダムテスト、障害率、信頼性評価、負荷分散、機械学習、A/Bテストでは確率が必要である。確率は単なる予想ではなく、標本空間、事象、分布を定めた上で扱う数学的な記述である。
\end{infobox}
\begin{todobox}
サイコロの離散確率分布を入れる。

画像生成用プロンプト

白背景、モノクロ、A4本文幅の棒グラフを作成する。横軸に1から6のサイコロの目、縦軸に確率を置き、すべての棒を同じ高さ1/6にする。右側に「確率の総和は1」「平均は3.5」と注記する。日本語ラベル、講義ノート風。
\end{todobox}
\section{10 数値精度・正確度・誤差}
計算機は有限個のビットで数を表す。そのため、任意の大きさの数や任意の実数を完全には表せない。数値解析では、この制約の中で効率よく、十分に正確な計算を行う方法を考える。

固定小数点は、数を一定間隔で配置して表す方式である。処理速度が重要な場合に有利なことがある。浮動小数点は、非常に大きな値と非常に小さな値を広く扱えるが、丸め誤差を伴う。

\begin{longtable}{>{\raggedright\arraybackslash}p{0.26\linewidth} >{\raggedright\arraybackslash}p{0.68\linewidth}}
\toprule
用語 & 説明 \\
\midrule
\endfirsthead
\toprule
用語 & 説明 \\
\midrule
\endhead
正確度 & 測定値または計算値が真の値にどれだけ近いかである。 \\
精密度 & 同じ対象を複数回測った値が互いにどれだけ近いかである。 \\
絶対誤差 & 近似値と真値の差の絶対値である。 \\
相対誤差 & 絶対誤差を真値の大きさで割った比である。 \\
有効数字 & 信頼できる桁を表す。 \\
オーバーフロー & 計算結果が表現可能範囲を超えることである。 \\
丸め & 近い表現可能な値へ置き換えることである。 \\
\bottomrule
\end{longtable}
\begin{infobox}{実務での注意}
金額、医療、制御、物理シミュレーションでは、数値表現の選択が重要である。浮動小数点で0.1を正確に表せないこと、丸め順序で結果が変わること、オーバーフローが起こることを前提に設計する必要がある。
\end{infobox}
\begin{todobox}
正確度・精密度・誤差の概念図を入れる。

画像生成用プロンプト

白背景、モノクロ、A4本文幅の2段図を作成する。上段に的を描き、点が真値に近いが散らばる例、点がまとまるが真値からずれる例、点が真値近くにまとまる例を示す。下段に「近似値」「真値」「絶対誤差」「相対誤差」の関係を数直線で示す。日本語ラベル、講義ノート風。
\end{todobox}
\section{11 代数構造}
代数構造は、集合と演算を組み合わせ、その演算が満たす規則を調べる枠組みである。抽象的に見えるが、データ型、演算の合成、暗号、形式手法、関数型プログラミング、仕様記述で役立つ。

\subsection{11.1 群・半群・モノイド}
群は、集合と一つの二項演算からなる構造であり、閉包性、結合性、単位元、逆元を満たす。演算順序を入れ替えても結果が同じ群は可換群またはアーベル群である。

\begin{longtable}{>{\raggedright\arraybackslash}p{0.26\linewidth} >{\raggedright\arraybackslash}p{0.68\linewidth}}
\toprule
構造 & 満たす条件の概要 \\
\midrule
\endfirsthead
\toprule
構造 & 満たす条件の概要 \\
\midrule
\endhead
半群 & 閉包性と結合性を持つ集合と二項演算である。 \\
モノイド & 半群に単位元を加えた構造である。文字列連結と空文字、自然数の加算と0などが例である。 \\
群 & モノイドに加えて、各要素に逆元が存在する構造である。整数の加算は群である。 \\
部分群 & 大きな群の中に含まれ、同じ演算で群になる部分集合である。 \\
巡回群 & 一つの生成元からすべての要素を生成できる群である。 \\
\bottomrule
\end{longtable}
\begin{infobox}{ソフトウェアでの見方}
モノイドは実務でも現れやすい。ログ文字列の連結、リストの結合、集計値の合成、設定のマージなどは、「結合できる」「空の値がある」という性質を持つことが多い。分散処理で部分結果を安全に合成する考え方にもつながる。
\end{infobox}
\subsection{11.2 環・体}
環は、集合に二つの演算を定めた構造である。第一の演算でアーベル群になり、第二の演算が結合的で、第二の演算が第一の演算に対して分配的である。整数全体に通常の加算と乗算を入れた構造は環である。

体は、0を除く要素が第二の演算についてアーベル群になる環である。有理数に通常の加算と乗算を入れた構造は体である。

\begin{todobox}
代数構造の関係図を入れる。

画像生成用プロンプト

白背景、モノクロ、A4本文幅の階層・比較図を作成する。左から「半群」「モノイド」「群」を条件追加の矢印で示す。別枠で「環」と「体」を示し、環は加法群と乗法半群と分配法則からなることを短く注記する。日本語ラベル、講義ノート風。
\end{todobox}
\section{12 工学微積分}
微積分は、連続的な変化を扱う数学である。工学微積分では、極限、連続性、微分、積分、超越関数、ベクトル解析を、工学上の問題に適用する。

\begin{longtable}{>{\raggedright\arraybackslash}p{0.26\linewidth} >{\raggedright\arraybackslash}p{0.68\linewidth}}
\toprule
概念 & 説明と用途 \\
\midrule
\endfirsthead
\toprule
概念 & 説明と用途 \\
\midrule
\endhead
極限 & 入力がある値に近づくとき、関数値がどこへ近づくかを見る。 \\
連続性 & 小さな入力変化が小さな出力変化に対応する性質である。 \\
微分 & 変化率を扱う。最適化、速度、勾配、感度分析に使う。 \\
積分 & 面積、累積量、総量を扱う。負荷の累積、エネルギー、確率密度に使う。 \\
超越関数 & 指数関数、対数関数、三角関数など、代数方程式だけでは表せない関数である。 \\
ベクトル解析 & 三次元空間のベクトル場の微分・積分を扱う。物理、制御、画像処理に関係する。 \\
\bottomrule
\end{longtable}
\begin{infobox}{ソフトウェアでの見方}
微積分は、機械学習の最適化、物理シミュレーション、制御、画像処理、信号処理、性能の傾向分析で使われる。初学者は、公式暗記よりも「変化率」「累積量」「近似」という意味を理解することが重要である。
\end{infobox}
\begin{todobox}
微分と積分の用途図を入れる。

画像生成用プロンプト

白背景、モノクロ、A4本文幅の左右比較図を作成する。左に曲線と接線を描き「微分＝瞬間の変化率」と示す。右に曲線下の面積を塗り「積分＝累積量」と示す。下部に「最適化」「シミュレーション」「制御」「データ分析」を用途として並べる。日本語ラベル、講義ノート風。
\end{todobox}
\section{13 新しい進展}
第17章は、数学的基礎が新しい応用分野にも広がっていることを示す。計算神経科学とゲノミクスは、数学モデル、シミュレーション、確率、グラフ、最適化、データ処理を必要とする分野であり、ソフトウェア工学とも深く関係する。

\subsection{13.1 計算神経科学}
計算神経科学は、神経系の働きを数学モデル、コンピュータシミュレーション、脳の抽象化によって理解しようとする分野である。脳を、情報を符号化し、時間とともに状態が変化する動的システムとして見る。

主な観点には、神経活動が情報をどう符号化するかを扱う神経符号化、単一ニューロンの生物物理、ニューロン同士の相互作用から生じるネットワークの動力学がある。機械学習、人工知能、制御理論、サイバネティクスとも接点がある。

\subsection{13.2 ゲノミクス}
ゲノミクスは、ゲノムの構造、機能、対応付け、進化、編集を扱う分野である。ゲノムは、DNAにより構成される遺伝情報であり、コンピュータ上で塩基配列を解析することをインシリコ解析と呼ぶ。

ゲノム解析では、DNA配列決定、生命情報学的解析、遺伝的変異の解析、計算モデル化などが使われる。大量の塩基配列データを扱うため、データ管理、セキュリティ、共有、解析効率が重要であり、専用ソフトウェアの開発にはソフトウェア工学が必要である。

\begin{todobox}
数学的基礎と新しい応用分野の関連図を入れる。

画像生成用プロンプト

白背景、モノクロ、A4本文幅の関連図を作成する。中央に「数学的基礎」を置き、左に「計算神経科学」、右に「ゲノミクス」を配置する。数学的基礎から「モデル化」「シミュレーション」「確率」「グラフ」「最適化」「大規模データ処理」の矢印を伸ばす。日本語ラベル、講義ノート風。
\end{todobox}
\section{14 トピックと参考文献の対応}
この表は、SWEBOKが示すトピックと参考資料の対応を、学習用に簡略化したものである。離散数学の基礎にはRosenを、数値計算にはCheney and Kincaidを読むとよい。

\begin{longtable}{>{\raggedright\arraybackslash}p{0.26\linewidth} >{\raggedright\arraybackslash}p{0.68\linewidth}}
\toprule
トピック & 主な参考資料 \\
\midrule
\endfirsthead
\toprule
トピック & 主な参考資料 \\
\midrule
\endhead
1. 基本論理 & Rosen, Discrete Mathematics and Its Applications \\
2. 証明技法 & Rosen \\
3. 集合・関係・関数 & Rosen \\
4. グラフと木 & Rosen \\
5. 有限状態機械 & Rosen \\
6. 文法 & Rosen \\
7. 数論 & Rosen \\
8. 数え上げの基礎 & Rosen \\
9. 離散確率 & Rosen \\
10. 数値精度・正確度・誤差 & Cheney and Kincaid, Numerical Mathematics and Computing \\
11. 代数構造 & Rosen \\
12. 工学微積分 & Cheney and Kincaid \\
\bottomrule
\end{longtable}
\section{15 発展的な読み物}
\begin{longtable}{>{\raggedright\arraybackslash}p{0.26\linewidth} >{\raggedright\arraybackslash}p{0.68\linewidth}}
\toprule
文献 & 読むと得られること \\
\midrule
\endfirsthead
\toprule
文献 & 読むと得られること \\
\midrule
\endhead
Rosen, Discrete Mathematics and Its Applications & 論理、証明、集合、関係、関数、グラフ、木、数論、数え上げ、離散確率を体系的に学べる。ソフトウェアの基礎数学として最初に読む価値が高い。 \\
Cheney and Kincaid, Numerical Mathematics and Computing & 数値計算、近似、誤差、数値アルゴリズムを学べる。シミュレーション、最適化、科学技術計算に関わる場合に有用である。 \\
\bottomrule
\end{longtable}
\section{16 参考文献}
1. K. Rosen, Discrete Mathematics and Its Applications, 8th ed., McGraw-Hill, 2018.

2. E. W. Cheney and D. R. Kincaid, Numerical Mathematics and Computing, 7th ed., Addison Wesley, 2020.

\section{17 画像TODO一覧}
本文に配置した画像TODOを再掲する。実際に図を作る場合は、各TODOのプロンプトをそのまま画像生成に使うと、A4資料に合う図を作りやすい。

\begin{itemize}[leftmargin=2zw]
\item 17章の全体地図。
\item 命題論理の真理値表と論理演算の図。
\item 量化子の意味を示す図。
\item 証明技法の比較フロー。
\item 集合演算のベン図。
\item 関係と関数の違いを示す図。
\item グラフの種類を比較する図。
\item 木構造と二分探索木の巡回を示す図。
\item 有限状態機械の状態遷移図と状態表。
\item 文法とチョムスキー階層の図。
\item 数の種類の包含図。
\item 和の規則・積の規則・順列・組合せの比較図。
\item サイコロの離散確率分布。
\item 正確度・精密度・誤差の概念図。
\item 代数構造の関係図。
\item 微分と積分の用途図。
\item 数学的基礎と新しい応用分野の関連図。
\end{itemize}
\section{18 章末確認問題}
\begin{itemize}[leftmargin=2zw]
\item 命題と述語の違いを、例を挙げて説明せよ。
\item 「すべての利用者は利用者IDを持つ」を述語論理の考え方で説明せよ。
\item 背理法と対偶証明は、どのような場面で使いやすいか。
\item 集合の共通部分、和集合、差集合を、権限管理の例で説明せよ。
\item 関係と関数の違いを説明せよ。API設計でこの違いが重要になる例を述べよ。
\item 有向グラフと重み付きグラフが、ソフトウェア開発で使われる例を一つずつ挙げよ。
\item 二分探索木を中順に巡回すると、なぜ昇順の列が得られるかを説明せよ。
\item 有限状態機械が画面遷移や通信プロトコルの設計に向く理由を説明せよ。
\item 正規文法、文脈自由文法、正規表現の関係を説明せよ。
\item 剰余演算がソフトウェアで使われる例を三つ挙げよ。
\item 全組合せテストが現実的でない場合、数え上げの知識がどう役立つかを説明せよ。
\item 離散確率分布において、各確率の合計が1でなければならない理由を説明せよ。
\item 正確度と精密度の違いを説明せよ。
\item モノイドが分散処理や集計に役立つ理由を説明せよ。
\item 微分と積分が機械学習やシミュレーションで使われる理由を説明せよ。
\end{itemize}
\begin{infobox}{解答の方向性}
1は真偽が決まる文と、変数を含む性質・関係の違いである。4は集合演算を利用者集合や権限集合に対応させる。8は「状態」と「入力」と「遷移」で振る舞いを表にできる点を述べる。13は、真値への近さと測定値同士の近さの違いを述べる。14は、結合演算と単位元があるため部分結果を安全に合成しやすい点を述べる。
\end{infobox}
\end{document}
